% from ega4-18-fr.tex
\oldpage[IV-4]{201}

il suffit de considérer le cas où $S$ est le spectre d'un corps algébriquement clos. Mais alors il revient au même de dire qu'un $S$-préschéma localement de type fini est lisse sur $S$ en un point ou qu'il est régulier en ce point (17.15.1), donc les conclusions résultent de (19.4.4).

Lorsque $X\to S$ est localement de présentation finie, et que l'immersion canonique $Y\to X$ est quasi-régulière, cette immersion est par définition de présentation finie puisque l'Idéal qui définit $Y$ est de type fini, donc le morphisme composé $Y\to X\to S$ est localement de présentation finie. On a déjà vu (19.4.3) que dans ces conditions le morphisme $Y'\to Y$ est lisse; en outre, on sait alors que $\sh{N}_{Y/X}$ est un $\sh{O}_Y$-Module localement libre et que $\operatorname{gr}_{\mathscr J}^*(\sh{O}_X)$ est isomorphe à $\mathbf S_{\sh{O}_Y}^*(\sh{N}_{Y/X})$, donc est un $\sh{O}_Y$-Module plat. Les hypothèses de (19.4.7) sont donc toutes vérifiées et l'on peut appliquer les conclusions de l'énoncé de (19.4.8).

\begin{corollary}[19.4.9]
\label{IV.19.4.9}
Sous les hypothèses générales de (19.4.8), supposons en outre que $Y$ soit lisse sur $S$ au point $x$. Alors, pour toute générisation $x_1$ de $x$ dans $X$ n'appartenant pas à $Y$, $X$ est lisse sur $S$ au point $x_1$. Si $Y$ est lisse sur $S$, alors il existe un voisinage ouvert $U$ de $Y$ dans $X$ tel que $X$ soit lisse sur $S$ aux points de $U-Y$.
\end{corollary}

Pour la première assertion, on peut se borner au cas où $X$ est un schéma local de point fermé $x$, et l'hypothèse entraîne alors que $Y$ est lisse sur $S$, tout voisinage du point fermé de $Y$ dans $Y$ étant nécessairement $Y$ lui-même; on conclut donc par le même raisonnement que dans (19.4.5), utilisant le fait que l'ensemble des points où $X'$ est lisse sur $S$ est ouvert dans $X'$. Pour prouver la seconde assertion, on peut se ramener au cas où $S$ est noethérien, et $X$ de type fini sur $S$, grâce à (8.9.1), (8.6.3), (11.2.9) et (17.7.6); on conclut alors comme dans (19.4.5), en utilisant le fait que l'ensemble des points où $X$ est lisse sur $S$ est ouvert dans $X$.

\begin{remark}[19.4.10]
\label{IV.19.4.10}
Dans (19.4.6, (ii)), remplaçons l'hypothèse que $X$ et $Y$ sont localement de présentation finie sur $S$ par celle que $S$ et $X$ (donc aussi $Y$) sont \emph{localement noethériens}. Alors il est clair que $\bigoplus_{n\geq0}\mathscr J^n$ est une $\sh{O}_X$-Algèbre de type fini, que le morphisme $X'\to X$ est de type fini, et il résulte encore de (19.2.4) appliqué aux préschémas localement noethériens, que l'immersion $Y'\to X'$ est transversalement régulière relativement à $S$ et que $X'$ est $S$-plat aux points de $Y'$.
\end{remark}

\begin{proposition}[19.4.11]
\EGAReaderTarget{ega4.c04.prop.19-4-11}{IV.19.4.11}
Soient $X$ un préschéma, $\mathscr E$ un $\sh{O}_X$-Module quasi-cohérent, $u:\mathscr E\to\sh{O}_X$ un homomorphisme de $\sh{O}_X$-Modules, $\mathscr J=u(\mathscr E)$ l'Idéal quasi-cohérent de $\sh{O}_X$ image de $u$, $Y$ le sous-préschéma fermé de $X$ défini par $\mathscr J$. Soit $X'$ le $X$-schéma obtenu en faisant éclater $Y$. D'autre part, soient $P=\mathbf P(\mathscr E)$ le fibré projectif sur $X$ défini par $\mathscr E$ (II, 4.1.1), $p:P\to X$ le morphisme structural, $\alpha_1^{\#}:p^*(\mathscr E)\to\sh{O}_P(1)$ l'homomorphisme canonique (II, 4.1.5.1) et $\mathscr H$ son noyau, de sorte qu'on a la suite exacte
\[
  0\longrightarrow\mathscr H\longrightarrow p^*(\mathscr E)\xrightarrow{\alpha_1^{\#}}\sh{O}_P(1)\longrightarrow0.
\]
Enfin, soit $\mathscr K$ l'Idéal quasi-cohérent de $\sh{O}_P$ image de la restriction $v:\mathscr H\to\sh{O}_P$ de $p^*(u)$, et soit $Z$ le sous-préschéma fermé de $P$ défini par $\mathscr K$.

